% Fichier généré automatiquement par tools/generate_corrections.py.
% Source : DYN/DYN-05-Methode/STOCK/04_RR_02/04_RR_02.tex
% Statut : complété (3/3 question(s)).
\begin{corrigebox}[Corrigé complété]
\CorrigeQuestion{1}
On choisit les inconnues et les conventions positives de la figure,
        on écrit les relations de conservation/fermeture adaptées, puis on résout le
        système obtenu. Le résultat symbolique doit être homogène et vérifié dans les cas
        limites (entrée nulle, rapport unitaire ou position de référence).
\CorrigeQuestion{2}
\CorrigeSource{Réponse reprise d'un exercice homonyme de la banque ; la provenance exacte est indiquée dans le commentaire du fichier.}
% Réponse reprise de : CIN/CIN-02-VitesseAcceleration/04_RR_02/04_RR_02.tex
$\torseurcin{V}{2}{0}=\torseurcin{V}{2}{1}+\torseurcin{V}{1}{0}$. Pour sommer les torseurs, il faut écrire les vecteurs vitesses au même point, ici en $C$. 

$\torseurcin{V}{2}{0}=\torseurl{ \left( \dot{\theta}+\dot{\varphi}\right)\vk{0}}{R\dot{\theta}\vj{1}+L\left( \dot{\theta}+\dot{\varphi}\right)\vj{2}}{C}$
\CorrigeQuestion{3}
\CorrigeSource{Réponse reprise d'un exercice homonyme de la banque ; la provenance exacte est indiquée dans le commentaire du fichier.}
% Réponse reprise de : CIN/CIN-02-VitesseAcceleration/04_RR_02/04_RR_02.tex
$\vectg{C}{2}{0}$ $ = \deriv{\vectv{C}{2}{0}}{\rep{0}}$ $= \deriv{R\dot{\theta}\vj{1}+L\left( \dot{\theta}+\dot{\varphi}\right)\vj{2}}{\rep{0}}$.

De plus,  $\deriv{\vj{1}}{\rep{0}} = \deriv{\vj{1}}{\rep{1}} + \vecto{1}{0}\wedge \vj{1}$
$ =\dot{\theta}\vk{0} \wedge \vj{1} $ $=-\dot{\theta}\vi{1}$ et 
$\deriv{\vect{j_2}}{\rep{0}} = \deriv{\vect{j_2}}{\rep{2}} + \vecto{2}{0}\wedge \vj{2}$
$ = \left( \dot{\theta}+\dot{\varphi}\right)\vk{0}\wedge \vj{2}$ $ = -\left( \dot{\theta}+\dot{\varphi}\right)\vi{2}$.

On a donc $\vectg{C}{2}{0}$ $=R\ddot{\theta}\vj{1}-R\dot{\theta}^2\vi{1}+L\left( \ddot{\theta}+\ddot{\varphi}\right)\vj{2}-L\left( \dot{\theta}+\dot{\varphi}\right)^2\vi{2}$.
\end{corrigebox}
